\Needspace{6\baselineskip}\Article{351}
Keeping all the notation of the preceding article, it is clear that each coefficient of the equation whose roots are the \(\beta\) sums \((\gamma,\lambda),\allowbreak (\gamma,\lambda^{\prime}),\allowbreak (\gamma,\lambda^{\prime\prime})\), etc., can be reduced to a form such as \[A+a(\beta\gamma,1)+a^{\prime}(\beta\gamma,g)\ldots+a^\varepsilon(\beta\gamma,g^{\alpha-1})\] and that the numbers \(A,\allowbreak a\), etc., are all integers. The equation whose roots are the \(\beta\) periods of \(\gamma\) terms contained in another period \((\beta\gamma,k\lambda)\) is derived from it by replacing every period \((\beta\gamma,\mu)\) throughout the coefficients by \((\beta\gamma,k\mu)\). Thus, if \(\alpha=1\), all the \(\beta\) periods of \(\gamma\) terms will be determined by an equation of degree \(\beta\), each of whose coefficients reduces to the form \(A+a(\beta\gamma,1)\), \textit{and is therefore a known quantity}, since \((\beta\gamma,1)=(n-1,1)=-1\). If, however, \(\alpha>1\), the coefficients of the equation whose roots are all the periods of \(\gamma\) terms contained in any given period of \(\beta\gamma\) terms will be known quantities as soon as the numerical values of all the \(\alpha\) periods of \(\beta\gamma\) terms are known. — The coefficients of these equations can often be calculated more conveniently, especially when \(\beta\) is not very small, by first finding the sums of powers of the roots and then deducing the coefficients from them by Newton's theorem, just as above in Article 349.

\textit{Example I.} For \(n=19\), we seek the equation whose roots are the sums \((6,1),\allowbreak (6,2),\allowbreak (6,4)\). Denoting these roots respectively by \(p,\allowbreak p^{\prime},\allowbreak p^{\prime\prime}\), and the required equation by \[x^3-Axx+Bx-C=0\] we obtain \[A=p+p^{\prime}+p^{\prime\prime},\quad B=pp^{\prime}+pp^{\prime\prime}+p^{\prime}p^{\prime\prime},\quad C=pp^{\prime}p^{\prime\prime}\] Hence \[A=(18,1)=-1\] Further, we have
\SourceLeaf{436}
\[\begin{gathered}
pp^{\prime}=p+2p^{\prime}+3p^{\prime\prime},\quad pp^{\prime\prime}=2p+3p^{\prime}+p^{\prime\prime},\\
p^{\prime}p^{\prime\prime}=3p+p^{\prime}+2p^{\prime\prime}
\end{gathered}\] whence \[B=6(p+p^{\prime}+p^{\prime\prime})=6(18,1)=-6\] Finally, \[\begin{aligned}
C=(p+2p^{\prime}+3p^{\prime\prime})p^{\prime\prime}&=3(6,0)+11(p+p^{\prime}+p^{\prime\prime})\\
&=18-11=7
\end{aligned}\] so the required equation is \[x^3+xx-6x-7=0\] Using the other method, we have \[p+p^{\prime}+p^{\prime\prime}=-1\] \[\begin{gathered}
pp=6+2p+p^{\prime}+2p^{\prime\prime},\quad p^{\prime}p^{\prime}=6+2p^{\prime}+p^{\prime\prime}+2p,\\
p^{\prime\prime}p^{\prime\prime}=6+2p^{\prime\prime}+p+2p^{\prime}
\end{gathered}\] whence \[pp+p^{\prime}p^{\prime}+p^{\prime\prime}p^{\prime\prime}=18+5(p+p^{\prime}+p^{\prime\prime})=13\] and similarly \[p^3+p^{\prime3}+p^{\prime\prime3}=36+34(p+p^{\prime}+p^{\prime\prime})=2\] From these, Newton's theorem gives the same equation as before.

II. For \(n=19\), we seek the equation whose roots are the sums \((2,1),\allowbreak (2,7),\allowbreak (2,8)\). Denoting these respectively by \(q,\allowbreak q^{\prime},\allowbreak q^{\prime\prime}\), we find \[\begin{gathered}
q+q^{\prime}+q^{\prime\prime}=(6,1),\\
qq^{\prime}+qq^{\prime\prime}+q^{\prime}q^{\prime\prime}=(6,1)+(6,4),\quad qq^{\prime}q^{\prime\prime}=2+(6,2)
\end{gathered}\] Thus, keeping the notation of the preceding example, the required equation will be \[x^3-pxx+(p+p^{\prime\prime})x-2-p^{\prime}=0\] The equation whose roots are the sums \((2,2),\allowbreak (2,3),\allowbreak (2,5)\), contained in \((6,2)\), is deduced from the preceding one by replacing \(p,\allowbreak p^{\prime},\allowbreak p^{\prime\prime}\) respectively by \(p^{\prime},\allowbreak p^{\prime\prime},\allowbreak p\). Applying the same substitution again gives the equation whose roots are the sums \((2,4),\allowbreak (2,6)\). \((2,9)\) contained in \((6,4)\).

\Needspace{9\baselineskip}\begin{center}\textit{The solution of the equation \(X=0\) is built on the preceding investigations.}\end{center}


\Needspace{6\baselineskip}\Article{352}
The preceding theorems, together with their attached corollaries, contain the main points of the whole theory, and the method for finding the values of the roots \(\Omega\) can now be given briefly.
\SourceLeaf{437}

First choose a number \(g\) that is a primitive root modulo \(n\), and find the least residues of the powers of \(g\) up to \(g^{n-2}\) modulo \(n\). Resolve \(n-1\) into factors, and into prime factors if we wish to reduce the problem to equations of the lowest possible degree. Let these, in an entirely arbitrary order, be \(\alpha,\allowbreak \beta,\allowbreak \gamma\ldots\zeta\), and put \[\frac{n-1}{\alpha}=\beta\gamma\ldots\zeta=a,\qquad\frac{n-1}{\alpha\beta}=\gamma\ldots\zeta=b,\quad\text{etc.}\] Distribute all the roots \(\Omega\) into \(\alpha\) periods of \(a\) terms; each of these in turn into \(\beta\) periods of \(b\) terms; each of these again into \(\gamma\) periods, etc. By the preceding article, find the equation of degree \(\alpha\), \((A)\), whose roots are those \(\alpha\) sums of \(a\) terms; their values will therefore become known by solving this equation.

But a difficulty arises here, since it seems uncertain to which root of the equation \((A)\) each sum is to be set equal, namely which root should be denoted by \((a,1)\), which by \((a,g)\), etc. This can be resolved as follows. We may use \((a,1)\) to denote any root of equation \((A)\); for each of its roots is a sum of \(a\) roots from \(\Omega\), and it is entirely arbitrary which root from \(\Omega\) is denoted by \([1]\). We may therefore clearly assume that one of the roots of which any given root of equation \((A)\) consists is expressed by \([1]\), so that this root of equation \((A)\) becomes \((a,1)\). The root \([1]\), however, is not yet fully determined by this: it remains entirely arbitrary or undetermined which of the roots constituting \((a,1)\) we wish to take for \([1]\). But as soon as \((a,1)\) is determined, all the other sums of \(a\) terms can also be deduced rationally from it (Article 346). This also shows that only a single root need be found by solving the equation. — The following, less direct method can also be used for this purpose. Take a definite root for \([1]\), \textit{that is}, put \([1]=\cos\frac{kP}{n}+i\sin\frac{kP}{n}\), choosing the integer \(k\) at will, provided it is not divisible by \(n\). Then \([2],\allowbreak [3]\), etc., will also indicate definite roots, and therefore the sums \((a,1),\allowbreak (a,g)\), etc., will denote definite quantities. Once these have been calculated with a light hand from sine tables, namely with enough accuracy to decide which are greater and which smaller, no doubt can remain about the notation by which the individual roots of equation \((A)\) are to be distinguished.

When all the \(\alpha\) sums of \(a\) terms have been found in this way, find, by the preceding article, the equation \((B)\) of degree \(\beta\), whose roots are the \(\beta\) sums of \(b\) terms contained in \((a,1)\); all the coefficients of this equation will be known
\SourceLeaf{438}
quantities. Since it is still arbitrary which of the \(a=\beta b\) roots contained in \((a,1)\) is denoted by \([1]\), any given root of equation \((B)\) can be expressed by \((b,1)\), for we may clearly assume that one of the \(b\) roots composing it is denoted by \([1]\). Thus find any one root of equation \((B)\) by solving it, set it \(=(b,1)\), and from it derive, by Article 346, all the other sums of \(b\) terms. This also gives us a check on the calculation, since those sums of \(b\) terms that belong to the same periods of \(a\) terms must always yield known sums. — In some cases it may be just as convenient to find \(\alpha-1\) other equations of degree \(\beta\), whose roots are respectively the individual \(\beta\) sums of \(b\) terms contained in the remaining periods of \(a\) terms, \((a,g),\allowbreak (a,gg)\), etc., and to find \textit{all} the roots both of these equations and of equation \(B\) by solving them. We must then decide, as above with the help of a sine table, to which periods of \(b\) terms the individual roots obtained in this way are to be set equal. Various other devices can also be used to make this decision, which cannot be fully explained here; one, however, for the case \(\beta=2\), is particularly useful and can be explained more briefly by examples than by rules, and will be found in the following examples.

After the values of all the \(\alpha\beta\) sums of \(b\) terms have been found in this way, equations of degree \(\gamma\) can in exactly the same way determine from them all the \(\alpha\beta\gamma\) sums of \(c\) terms. Namely, we must find \textit{either one} equation of degree \(\gamma\), whose roots are the \(\gamma\) sums of \(c\) terms contained in \((b,1)\), by Article 350; obtain any one root by solving it, set it \(=(c,1)\), and finally deduce from it, by Article 346, all the other similar sums; \textit{or} we must develop in the same way a total of \(\alpha\beta\) equations of degree \(\gamma\), whose roots are respectively the \(\gamma\) sums of \(c\) terms contained in the individual periods of \(b\) terms, extract the values of all the roots of all these equations by solving them, and finally determine their order as above with the help of a sine table, or, for \(\gamma=2\), by the device to be shown below in the examples.

Continuing in this way, we will clearly at last obtain all the \(\frac{n-1}{\zeta}\) sums of \(\zeta\) terms. Thus, if we develop by Article 348 the equation of degree \(\zeta\) whose roots are the \(\zeta\) roots from \(\Omega\) contained in \((\zeta,1)\), all its coefficients will be known quantities. If any one of its roots is obtained by solving it, we may set it \(=[1]\), and all the other roots \(\Omega\) will be obtained from its powers.
\SourceLeaf{439}
If preferred, all the roots of that equation can also be found by solving it, and, in addition, by solving \(\frac{n-1}{\zeta}-1\) other equations of degree \(\zeta\), which respectively give all the \(\zeta\) roots contained in each of the remaining periods of \(\zeta\) terms, all the other roots \(\Omega\) can be found.

It is also clear that, as soon as the first equation \((A)\) is solved, or as soon as the values of all the \(\alpha\) sums of \(a\) terms are known, Article 348 immediately gives the resolution of the function \(X\) into \(\alpha\) factors of degree \(a\). Further, after solving equation \((B)\), or after finding the values of all the \(\alpha\beta\) sums of \(b\) terms, each of those factors is again resolved into \(\beta\), or \(X\) into \(\alpha\beta\) factors of degree \(b\), etc.


